\pdfsuppresswarningpagegroup=1
\documentclass[mytheorems]{SIAMbook2019}

\usepackage{graphicx}
\usepackage{multicol}
\usepackage{amsmath}
\usepackage[nottoc]{tocbibind}
\usepackage{xcolor}
\usepackage{amsmath}
\usepackage{amssymb}
\usepackage{array}
\usepackage{booktabs}
\usepackage{tabularx}
\usepackage{mathtools}
\usepackage{enumitem}
\usepackage{multirow}
\usepackage{cancel}
\usepackage{fix-cm}
\usepackage{siunitx}

\usepackage{amsfonts}
\usepackage{times}
\usepackage{esint}
\usepackage[T1]{fontenc}
\renewcommand{\sfdefault}{phv}
\renewcommand{\ttdefault}{txtt}
\renewcommand\bfdefault{b}

\definecolor{tab-blue}{HTML}{1f77b4}
\definecolor{tab-orange}{HTML}{ff7f0e}
\definecolor{tab-green}{HTML}{2ca02c}
\definecolor{tab-red}{HTML}{d62728}
\definecolor{tab-purple}{HTML}{9467bd}
\definecolor{tab-brown}{HTML}{8c564b}
\definecolor{tab-pink}{HTML}{e377c2}
\definecolor{tab-gray}{HTML}{7f7f7f}
\definecolor{tab-olive}{HTML}{bcbd22}
\definecolor{tab-cyan}{HTML}{17becf}

\usepackage[backref=section,
            bookmarks=true,
            colorlinks=true,
            citecolor=tab-red,
            linkcolor=tab-blue,
            urlcolor=tab-blue,
            hypertexnames=true,
            linktoc=all]{hyperref}

\usepackage{latexsym}
\renewcommand{\sfdefault}{phv}
\renewcommand{\ttdefault}{txtt}
\renewcommand\bfdefault{b}

\usepackage{bm}
\renewcommand{\vec}[1]{\boldsymbol{#1}}
\newcommand{\dvec}[1]{\vec{#1}}

\definecolor{latexgray}{gray}{0.95}
\newcommand{\erratafix}[1]{\medskip\newline\colorbox{latexgray}{\begin{minipage}[t]{0.90\textwidth}#1\end{minipage}}}

% ----- algorithms
\usepackage[algo2e,ruled,vlined,algochapter,dotocloa,linesnumbered]{algorithm2e}
\SetCommentSty{scriptsize}
\SetKwComment{tcc}{\{}{\}}
\SetKwIF{If}{ElseIf}{Else}{if}{}{else if}{else}{endif}
\SetKwFor{While}{while}{}{endw}
\SetKwFor{ForEach}{for each}{}{}
\SetKwFor{For}{for}{}{}
\SetKwInOut{Input}{Input~}
\SetKwInOut{Output}{Output~}
\SetKwInput{Input}{Input~}
\SetKwInput{Output}{Output~}
\SetKw{Continue}{continue}
\SetKw{Break}{break}
\let\oldnl\nl% Store \nl in \oldnl
\newcommand{\nonl}{\renewcommand{\nl}{\let\nl\oldnl}}% Remove line number for one line of algorithm

\begin{document}

\thispagestyle{empty}
\par\noindent\hfill\rule{\textwidth}{0.5pt}

\medskip
\noindent{\Large Numerical Partial Differential Equations}\newline
\noindent James H. Adler, Hans De Sterck, Scott MacLachlan, Luke Olson
\medskip

\noindent Errata for the first printing. \today

\medskip
\noindent \url{https://epubs.siam.org/doi/book/10.1137/1.9781611978285}

\par\noindent\hfill\rule{\textwidth}{0.5pt}

\medskip

\begin{itemize}[itemsep=10pt]
  \item \textcolor{tab-red}{page 222, exercise \#7:} The Fromm slope should have a minus (not a plus):
    \erratafix{
    with $a > 0$.  Using the REA algorithm (see~Algorithm 6.1), consider the Fromm slope
      \begin{equation*}
      \delta_{k,\ell} = \frac{1}{2}(u_{k+1,\ell} - u_{k-1,\ell}).
      \end{equation*}
      }
  \item \textcolor{tab-red}{page 289, Algorithm 7.14: CG applied to $L^T A L\vec{y} = L^T\vec{f}$:, Line 7}
    The the final line should not have $L^T A L$.
    \erratafix{
      {\renewcommand{\thealgocf}{7.14}
      \begin{algorithm2e}[H]
  \caption{CG applied to $L^T A L\vec{y} = L^T\vec{f}$}
  \DontPrintSemicolon%
  \Input{$A$, linear system matrix\newline
         $L$, factorized preconditioner with $M = LL^T$\newline
         $\vec{f}$, linear system right-hand side\newline
         $\vec{u}_0$, initial guess
       }
   \Output{$\hat{\vec{y}}_{k+1}$, $(k+1)\text{th}$ iterate with $\vec{u}\approx L \hat{\vec{y}}_{k+1}$}
  \nonl\;
  $\hat{\vec{r}}_0 = L^T(\vec{f} - A\vec{u}_0)$, $\hat{\vec{p}}_0 = \hat{\vec{r}}_0$, $\hat{\vec{y}}_0 = L^{-1}\vec{u}_0$\;
  \For{$k=1,2, \ldots$}{
        $\alpha_k = \hat{\vec{r}}_k^T\hat{\vec{r}}_k / \hat{\vec{p}}_k^T L^T AL\hat{\vec{p}}_k$\;
        $\hat{\vec{y}}_{k+1} = \hat{\vec{y}}_{k} +\alpha_{k}\hat{\vec{p}}_k$\;
        $\hat{\vec{r}}_{k+1} = \hat{\vec{r}}_{k} -\alpha_{k}L^T A L\hat{\vec{p}}_k$\;
        $\beta_k = \hat{\vec{r}}_{k+1}^T\hat{\vec{r}}_{k+1} / \hat{\vec{r}}_{k}^T\hat{\vec{r}}_{k}$\;
        $\hat{\vec{p}}_{k+1} = \hat{\vec{r}}_{k+1} +\beta_{k} \hat{\vec{p}}_k$\;
  }
      \end{algorithm2e}}
      }
  \item  \textcolor{tab-red}{page 401, last paragraph:} The notation here is inconsistent, introducting $\phi_{q,k}(x)$, which later becomes $\phi_{k,q}(x)$.  The latter notation is more consistent:
 \erratafix{
For $v(x) \in \mathcal{V}^{h_x}$, the restriction to each element, $v_k(x)$, is a
polynomial of degree at most $p$.  Here, we consider a \textit{nodal} representation of the
polynomial, writing $v_k(x)$ in element $\Omega_k$ in terms of a nodal basis, $\{\phi_{k,q}(x)\}_{q=0}^p$,
where the basis functions, $\phi_{k,q}(x)$, satisfy $\phi_{k,q}(x_{k,j})=\delta_{q,j}$ at the
$p+1$ interpolation points, $x_{k,0},\dots,x_{k,p}$, in element $\Omega_k$, as shown in Figure 9.5.2.
The nodal basis functions $\phi_{k,q}(x)$ are given by
the Lagrange polynomials.
}

  \item  \textcolor{tab-red}{page 402, line 9:} Missing $\;\widehat{}\;$ on $\phi_q$:
    \erratafix{
    the nodes $x_{k,q}\in\Omega_k$, $\phi_{k,q}(x)$, are related as $\phi_{k,q}(F_k(\widehat{x})) = \widehat{\phi}_q(\widehat{x})$.
  }

  \item  \textcolor{tab-red}{page 402, two lines after (9.157):} The function $u^{h_x}$ should just be $u$:
    \erratafix{
equivalent, since we use the parts of $u$ and $v$ that are restricted to $\Omega_k$, and are,\ldots
}

  \item  \textcolor{tab-red}{page 404, eqn. (9.165):} The function $\vec{z}$ should depend on $\vec{u}_{k-1}$ and $\vec{u}_{k+1}$ in addition to $\vec{u}_k$:
    \erratafix{
\begin{align*}
  M \frac{d\vec{u}_k}{dt} & = \frac{2}{h_x} \left(K \vec{f}_k(\vec{u}_k)
  +f^*(u_{k-1,p},u_{k,0})
   \vec{\delta}_0
  -f^*(u_{k,p},u_{k+1,0})
  \vec{\delta}_p\right),\label{eq:dgsemi}\\
                      & =: -\vec{z}(\vec{u}_{k-1},\vec{u}_k,\vec{u}_{k+1}),
\end{align*}
}

\item  \textcolor{tab-red}{page 410, eqn. (9.187):} The sign on the exponent should be positive, not negative:
  \erratafix{
  \begin{equation*}
  \widetilde{f}_{q,q} =
  \begin{cases}
    1 & \text{for $q < p_{\text{cutoff}}$,}\\
    e^{\alpha\left(\frac{q+1-p_{\text{cutoff}}}{p+1-p_{\text{cutoff}}}\right)^{2s}} & \text{for $q\geq p_{\text{cutoff}}$,}
  \end{cases}
\end{equation*}
}

\item \textcolor{tab-red}{page 563, ref. [2]:} The term ``editors'' should be removed from the citation:
    \erratafix{
      {\sc R.~A. Adams and J.~J.~F. Fournier}, {\em Sobolev Spaces}, vol.~140 of Pure and Applied Mathematics, Elsevier/Academic Press, second~ed., 2003, \url{https://www.sciencedirect.com/bookseries/pure-and-applied-mathematics/vol/140}.
    }

\end{itemize}

\end{document}
